Báo cáo phân tích \href{https://arxiv.org/abs/2206.02066}{\emph{PIDNet: A Real-Time Semantic Segmentation Network Inspired by PID Controllers}} của Jiacong Xu, Zixiang Xiong và Shankar P. Bhattacharyya, công bố tại CVPR 2023 \cite{pidnet}. Paper đề xuất kiến trúc ba nhánh để xử lý việc đặc trưng ngữ cảnh lấn át chi tiết ở biên vật thể: P giữ chi tiết, I tổng hợp ngữ cảnh và D học thông tin biên để điều tiết fusion. Phần phân tích làm rõ cơ sở PID, vai trò của Pag, PAPPM, Bag và các loss, đồng thời xem xét bằng chứng ablation và giới hạn của kết quả công bố.

Thực nghiệm khảo sát quy tắc boundary-guided fusion trên Oxford-IIIT Pet bằng hai biến thể ResNet18 có cùng boundary supervision, lần lượt dùng hệ số cố định và hệ số phụ thuộc vị trí. Qua ba seed, mỗi lượt 12 epoch trên toàn bộ tập train, mIoU đạt $78{,}00\pm0{,}11$\% với fusion cố định và $77{,}82\pm0{,}03$\% với fusion theo biên. Chênh lệch ghép theo seed là $-0{,}18\pm0{,}12$ điểm phần trăm. IoU dải chuyển tiếp lần lượt là $55{,}15\pm0{,}39$\% và $55{,}02\pm0{,}13$\%. Boundary head học được tín hiệu phân biệt các vùng, nhưng dùng trực tiếp tín hiệu này làm hệ số fusion chưa cải thiện mIoU. Kết quả được diễn giải trong phạm vi phép trộn được khảo sát và những khác biệt với PIDNet đầy đủ.
